\subsection{Iteration and the description of the final tables}
\label{iter:subsection}

Fix the round just constructed, including $D$, $G$, $q$ and its $s$
grouping steps. To save a factor $h^a$ in the total charge, take
\[
 j=\left\lceil\frac{a\log h}{\log G}\right\rceil,\qquad
 L_i=L_0G^{-i},\qquad T_i=T_0q^i \qquad(0\le i\le j).
\]
Here $L_i,T_i$ are the budgets after $i$ completed rounds.
It remains to choose $r$ and the initial number $S$ of blocks. A large
$r$ makes partially occupied packs inexpensive. A large $S$ supplies
enough groups for all $sj$ steps and leaves enough blocks to amortize
the common host in Section~\ref{sec:assembly}.

\paragraph{Packing slack and exterior width.}
In one grouping step $N_b$ decreases by $r$ while $H$ increases by at
most $cr$. Their product, which counts allocated occurrences, increases
by at most $c$. The fixed expansions of the host can be included in $c$.
Thus at every stage of the construction,
\begin{equation}
 H N_b\le C^{j+1}mS\log h , \label{iter:occurrences}
\end{equation}
where $C$ depends only on the design of one round. Since
$L_0=C_0 mT_0 S\log(2+h)$, the ratio that controls the slack during
the round starting from $L_i,T_i$ is bounded by
\[
   \frac{H N_b T_i}{rL_i}\le \frac{C_1C^{j+1}(qG)^i}{r}.
\]
Choose $r$ to be the least power of two at least $C_2(CqG)^{j+1}$.
The round design fixes $C_2$ so that this ratio is at most $c_D$ in
\eqref{eq:round-slack} and $r$ meets the lower bound in every step.
In particular $r=h^{O_D(a)+o(1)}$, independently of $S$. The external
width has a separate bound. The $ed$ exposed bits are present in every
round, and $G^j<Gh^a$. Hence for $i\le j$,
\[
       \frac{ed}{L_i}=\frac{eG^i}{C_0m\log(2+h)}
           =O_D\!\left(\frac{h^{-a}}{\log h}\right).
\]
For sufficiently large $h$ this gives $ed\le L_i/D$, the width
allowance for the immobile exterior identities.

\paragraph{Copies and terminal budgets.}
Let $R$ be the least power of two at least $h^M$, for any fixed $M>0$.
Taking $S=Rr^{sj}$ leaves exactly $R$ blocks after $j$ rounds.
Each initial block contains $T_0$ copies, so $d=T_0S$ and
$\log d=O((\log h)^2)$. In particular $d=n^{o(1)}$.
Now $h^a\le G^j<Gh^a$, $q\le G^{1+\delta/3}$ and $T_0=\Theta(h^b)$.
It follows from \eqref{eq:round-start} and \eqref{iter:occurrences} that
\begin{equation}
 \begin{aligned}
 L_j&=O(md\log h\,h^{-a}),&
 T_j&=O\bigl(h^{b+(1+\delta/3)a}\bigr),\\
 HN_b&\le mS h^{O_D(a)+o(1)} .
 \end{aligned}
 \label{eq:round-terminal}
\end{equation}
The occurrence bound holds throughout, with a constant in $O_D(a)$
independent of $M,a,b,\beta$. These scales also justify using a single
working field. They range from $T_0\ge4h^\beta$ to a fixed power of
$h$. Choose the extension degree divisible by the degree of the fixed
extension in Lemma~\ref{lem:local}, and large enough to meet the
cardinality requirement at $T_j$. Its degree is $O(\log h)$ and hence
$O(\log T_i)$ for every $i$.

\paragraph{Descriptions of tables.}
An occurrence carries a label for its table, the widths of its bundles,
and their positions in its port order. The cost of naming its
incident host edges is bounded in terms of $D$ because each edge can
be named by its place in the local neighbor list. Constants indexed by $D$ may depend on the round design. These local data have total length per occurrence at most
\begin{equation}
 C_3^{j+1}(\log h)^{O(1)}
   =h^{O_D(a)+o(1)} .
 \label{iter:labels}
\end{equation}

At initialization the stack at a position has a fixed table label,
specified by its original table, its stack length and the locations
of its ordered sublists. Its description has length
$(\log h)^{O(1)}$. Within one round the following facts bound
the number of old labels that can enter any new label.
\begin{itemize}
\item A colony of capacity $DT_i$ contains at most $D/\theta$ live
entering stations, since each is charged at least $\theta T_i$.
The local expansion in Lemma~\ref{lem:local} uses a bounded number of
old labels, discrete choices and field elements in each output
description. Its additional description length is $O_D(\log h)$.
\item In a subround or reset, both the bound on load per leaf and the
bound on the number of leaves contributing to one resulting pack are
fixed multiples, depending on $D$, of $T_i$ and $1$, respectively.
Each live core contributes at least one internal unit of width
$u\ge T_i/\operatorname{poly}(D)$. The charged load of an individual
frozen station has the same kind of lower bound even after the fixed
rescaling. Scalar stations are withdrawn as soon as they form. It
follows that a contributing pack has a bounded number of old labels
and a bounded number of units. Lists of unit pairings, rotor layers
and ordered sublists at a port therefore have bounded numbers of
entries. Each entry needs at most $O_D(\log h+\log r)$ bits. This
also describes the routing on the final trees at the frozen stations.
\end{itemize}
Thus if $A$ bounds the current label length, one step gives a bound
$C_4A+C_4(\log h)^{C_5}$, including the descriptors of locally
withdrawn scalar factors. The multiplicative constant counts old
labels, independently of the field degree and the length of the input
dictionary labels. Iteration through the $sj$ steps proves
Equation~\eqref{iter:labels}. Global host addresses enter the control
of a whole position only. They never have to be appended to an
individual table description.

\paragraph{The tensor represented by the sum.}
Each term now consists of $R$ surviving blocks and the scalar factors
withdrawn along its history. Regrouping permutes the exterior lists along with their stations and preserves each product entry by entry. A local tensor identity replaces a product by a sum with the same entries on those $ed$ named ports. Colorings, packings, and routings use fixed tie-breaking rules. Only choices from local tensor identities index the expansion.
The final internal wire assignments are summed when the $R$ blocks are
contracted.

For this evaluation one dictionary suffices. A key records the type,
stage, description and arity of a table, followed by its incident bits.
Length fields or fixed padding give a unique encoding. The input
dictionary and the rules of the construction assign a field value
to each valid key. The values at unused keys may be fixed arbitrarily.
In particular, a key for a withdrawn scalar has no incident bits.
A longest key has length
\[
 O_D(T_j)+h^{O_D(a)+o(1)}
\]
by the maximum-arity invariant and Equation~\eqref{iter:labels}.
A derived table description can exceed $(\log h)^{O(1)}$ bits. Section~\ref{sec:assembly} places that description in the \emph{argument} of a dictionary call. The label of the installed table need only specify the key format and length. This accounts separately for the arity of a lookup and the short label required as input to the next macro.
